\providecommand{\Vir}{\mathrm{Vir}}
\providecommand{\Muk}{\mathrm{Muk}}
\providecommand{\sltwo}{\mathfrak{sl}_2}
\providecommand{\Heis}{\mathrm{Heis}}
\providecommand{\DS}{H^{0}_{\mathrm{DS}}}
\providecommand{\Eone}{E_1}
\providecommand{\Etwo}{E_2}
\providecommand{\EZ}{\mathrm{EZ}}
\providecommand{\Dfive}{D_5}
\providecommand{\PhiTen}{\Phi_{10}}
\providecommand{\PhiUn}{\Phi_{10}^{\mathrm{un}}}
\providecommand{\PhiOP}{\Phi_{10}^{\mathrm{OP}}}

\section{Fourier match through $q^{24}$: Heisenberg--Mukai and the Humbert residue}
\label{wn:sec:fourier-match-q24}

\paragraph{Normalisations.}
Let
\[
  Z = \begin{pmatrix}\tau & z\\ z & \sigma\end{pmatrix}\in\mathcal H_2,
  \qquad
  q=e^{2\pi i\tau},\quad p=e^{2\pi i\sigma},\quad
  \zeta=e^{2\pi i z}.
\]
The Eichler--Zagier generator of \(J_{0,1}^{\mathrm{weak}}\) is
\[
  \phi_{0,1}^{\EZ}(\tau,z)
  =
  \zeta^{-1}+10+\zeta
  +q(10\zeta^{-2}-64\zeta^{-1}+108-64\zeta+10\zeta^2)
  +O(q^2).
\]
Its discriminant-zero coefficient is \(c_1(0)=10\). The primitive
Borcherds denominator is
\[
  \Dfive=\operatorname{Borch}(\phi_{0,1}^{\EZ})
  =64^{-1}\Delta_5,\qquad
  \kappa_{\mathrm{BKM}}(\Delta_5)
  =\mathrm{wt}(\Dfive)=\frac{c_1(0)}{2}=5.
\]
The symbol \(\Delta_5\) denotes the unnormalised Gritsenko theta-product
representative, while \(\Dfive\) denotes the monic Borcherds product.
Doubling the Jacobi input doubles the Borcherds exponents:
\[
  \operatorname{Borch}_{\Delta}(2\phi_{0,1}^{\EZ})
  =\PhiUn=\Delta_5^2,\qquad
  \operatorname{Borch}_{D}(2\phi_{0,1}^{\EZ})
  =\PhiOP=\Dfive^2=4096^{-1}\Delta_5^2.
\]
The Oberdieck--Pandharipande reduced-DT scalar uses the second square:
\[
  Z^{\mathrm{red}}_{\mathrm{DT}}(K3\times E)
  =-(\PhiOP)^{-1}
  =-\Dfive^{-2}
  =-4096\,\Delta_5^{-2}.
\]
The Humbert-residue computation below uses the primitive inverse
\(\Delta_5^{-2}\), hence the OP/DT scalar differs by the external
factor \(-4096\).

\paragraph{Humbert residue.}
In the unnormalised Igusa square convention \(\PhiUn=\Delta_5^2\), the
separating degeneration is
\[
  \PhiUn(\tau,z,\sigma)
  =
  (2\pi i z)^2\eta(\tau)^{24}\eta(\sigma)^{24}+O(z^4)
\]
in this normalisation (Gritsenko--Nikulin 1996, Corollary~5.2;
Borcherds 1998, Theorem~10.1; van der Geer 2008, Lecture~6,
Equation~(6.12)). Hence on the diagonal Humbert surface
$H_1=\{\sigma=\tau\}$,
\[
  R(\tau)
  :=
  \bigl[(2\pi i z)^2\Delta_5^{-2}\bigr]_{H_1,z=0}
  =
  \eta(\tau)^{-48}
  =
  q^{-2}\prod_{n\geq 1}(1-q^n)^{-48}.
\]
The displayed degeneration fixes the sign: the quotient
\[
  \frac{(2\pi i z)^2\Delta_5^{-2}}
       {\eta(\tau)^{-24}\eta(\sigma)^{-24}}
\]
tends to \(+1\) as \(z\to0\).

\paragraph{Coefficient computation.}
Set
\[
  g(q):=q^2R(\tau)=\prod_{n\geq1}(1-q^n)^{-48}
       =\sum_{n\geq0}g_nq^n.
\]
The coefficients are the $48$-coloured partition numbers
\[
  g_n=\sum_{\lambda\vdash n}\prod_{i\geq1}
       \binom{47+m_i(\lambda)}{m_i(\lambda)},
\]
where $m_i(\lambda)$ is the multiplicity of the part $i$ in
$\lambda$. Equivalently, the logarithmic derivative gives the recurrence
\[
  g_0=1,\qquad
  n\,g_n =48\sum_{m=1}^{n}\sigma_1(m)\,g_{n-m}
  \quad (n\geq1).
\]
The truncated Euler product, the recurrence, and the partition sum give
the same integer vector through \(q^{24}\):
{\small
\begin{align*}
  &g_0 = 1,\quad g_1 = 48,\quad g_2 = 1{,}224,\quad g_3 = 21{,}952,\\
  &g_4 = 309{,}876,\quad g_5 = 3{,}657{,}312,\quad
    g_6 = 37{,}468{,}928,\\
  &g_7 = 341{,}773{,}440,\quad
    g_8 = 2{,}826{,}752{,}418,\\
  &g_9 = 21{,}491{,}641{,}808,\quad
    g_{10} = 151{,}810{,}235{,}136,\\
  &g_{11} = 1{,}004{,}753{,}937{,}600,\quad
    g_{12} = 6{,}273{,}891{,}838{,}360,\\
  &g_{13} = 37{,}171{,}410{,}206{,}112,\quad
    g_{14} = 209{,}969{,}121{,}051{,}648,\\
  &g_{15} = 1{,}135{,}389{,}617{,}917{,}568,\quad
    g_{16} = 5{,}897{,}908{,}848{,}093{,}087,\\
  &g_{17} = 29{,}521{,}227{,}582{,}821{,}520,\quad
    g_{18} = 142{,}760{,}699{,}405{,}228{,}800,\\
  &g_{19} = 668{,}557{,}347{,}682{,}912{,}128,\quad
    g_{20} = 3{,}038{,}278{,}325{,}284{,}762{,}308,\\
  &g_{21} = 13{,}423{,}984{,}754{,}534{,}844{,}000,\\
  &g_{22} = 57{,}759{,}538{,}443{,}660{,}433{,}920,\\
  &g_{23} = 242{,}385{,}485{,}861{,}920{,}278{,}528,\\
  &g_{24} = 993{,}392{,}557{,}953{,}227{,}803{,}294.
\end{align*}%
}
The finite table proves the displayed truncation. The all-degree product
identity comes from the separating degeneration together with the
rank-\(48\) Heisenberg character formula.

\paragraph{Heisenberg character.}
The Mukai lattice
\[
  \Muk(K3)=H^0(K3,\mathbb Z)\oplus H^2(K3,\mathbb Z)\oplus H^4(K3,\mathbb Z)
  \cong \mathrm{II}_{4,20}
\]
has rank $24$. The doubled Mukai Heisenberg algebra
\[
  \Heis(\Muk(K3)^{\oplus2})
\]
therefore has rank $48$. Its unshifted Fock character is
\[
  \chi_{\Heis}^{\mathrm{Fock}}(q)
  =
  \prod_{n\geq1}(1-q^n)^{-48}=g(q),
\]
and its conformally shifted VOA character is
\[
  q^{-48/24}\chi_{\Heis}^{\mathrm{Fock}}(q)
  =
  q^{-2}g(q)=\eta(\tau)^{-48}=R(\tau).
\]
Thus the coefficient comparison through $q^{24}$ is exactly the
rank-$48$ Heisenberg Fock computation after the Humbert residue
specialisation of $\Delta_5^{-2}$.

\paragraph{Shapovalov form.}
The preceding character belongs to the Heisenberg Fock module. For a
finite-rank Heisenberg algebra with non-degenerate form and non-zero
central level, the Fock module is
irreducible (Kac 1990, Chapter~2; Frenkel--Ben-Zvi 2004,
Chapter~3). On the degree-$N$ subspace the Shapovalov form decomposes
over partitions of $N$ into symmetric powers of the Mukai pairing,
with determinant factors coming only from the oscillator pairing. The
Virasoro Kac determinant
\[
  \prod_{rs\leq N}\bigl(h-h_{r,s}(c)\bigr)^{p(N-rs)}.
\]
belongs to Virasoro Verma modules and has no role in the rank-$48$
Heisenberg dimension count.

\paragraph{Central charges.}
Three central-charge computations occur in adjacent literature. They are
attached to three different objects.

\begin{enumerate}[leftmargin=1.5em]
\item The doubled Heisenberg VOA has conformal central charge $48$.
This accounts for the shift $q^{-48/24}=q^{-2}$ in
$\eta(\tau)^{-48}$.
\item The class-$\mathcal S$ theory
$\mathcal T[A_1,\Sigma_{0,24}]$ has
$(n_v,n_h)=(63,88)$ in the Chacaltana--Distler pants decomposition.
Hence
\[
  c_{4d}=\frac{2n_v+n_h}{12}=\frac{2\cdot63+88}{12}
        =\frac{107}{6},
  \qquad
  c_{2d}=-12c_{4d}=-214
\]
by the Beem--Rastelli $4d/2d$ dictionary. This is class-$\mathcal S$
Schur-sector data, separate from the Heisenberg central charge.
\item For the principal Drinfeld--Sokolov reduction of
$\widehat{\sltwo}$ at $k=-2+1/22$, the Virasoro central charge is
\[
  c(k)=1-\frac{6(k+1)^2}{k+2}
      =1-6\left(\frac{21}{22}\right)^2\cdot22
      =-\frac{1312}{11}.
\]
Twenty-two uncoupled copies give $-2624$. The class-$\mathcal S$ value
\(-214\) requires the full class-$\mathcal S$/BRST construction. The
tensor product of these Virasoro quotients gives the wrong central
charge.
\end{enumerate}

\paragraph{Virasoro minimal models.}
The Virasoro minimal model labelled by $(2,45)$ has
\[
  c_{2,45}
  =
  1-\frac{6(45-2)^2}{2\cdot45}
  =
  1-\frac{1849}{15}
  =
  -\frac{1834}{15}.
\]
This differs from $-214=-3210/15$ and from the Heisenberg value $48$.
Consequently the $(2,45)$ Kac table and its null-vector enumeration
belong to a separate coefficient problem from the integers $g_n$ above.

\paragraph{Borcherds constants.}
The primitive CHL denominator ladder carries
\[
\begin{array}{c|ccccc}
N & 1 & 2 & 3 & 4 & 6\\ \hline
c_N(0) & 10 & 8 & 6 & 4 & 2\\
\kappa_{\mathrm{BKM}}(\Phi_N) & 5 & 4 & 3 & 2 & 1
\end{array}
\]
with
\[
  \kappa_{\mathrm{BKM}}(\Phi_N)=\mathrm{wt}(\Phi_N)=\frac{c_N(0)}{2}
\]
by the Borcherds weight theorem (Borcherds 1998, Theorem~13.3) and the
Gritsenko--Nikulin CHL construction. At \(N=1\), \(\Phi_1=\Delta_5\)
and \(\kappa_{\mathrm{BKM}}(\Delta_5)=5\).

The Nikulin extension and the Gritsenko--Clery simplest-divisor
catalogue use different index sets. The Nikulin-admissible extension
adds \(N=5,7,8\) with weights \((2,1,1)\), giving the sorted
twice-weight tuple
\[
  (10,8,6,4,4,2,2,2)\qquad (N=1,\ldots,8).
\]
The Gritsenko--Clery catalogue is indexed by triples
\[
  (t,N;k)\in
  \bigl\{(1,1;5),(2,1;2),(1,2;3),(3,1;1),
  (1,3;2),(4,1;\tfrac12),(1,4;\tfrac32),(2,2;1)\bigr\},
\]
whose twice-weight tuple is
\[
  (10,4,6,2,4,1,3,2).
\]
The shared row is \(F_{1,1}=\Delta_5\). All other rows carry their own
level, cover, multiplier, and denominator data.

\paragraph{\texorpdfstring{$K3\times E$}{K3 x E} invariants.}
The compact categorical Euler invariant of the total space is
\[
  \kappa_{\mathrm{cat}}(K3\times E)
  =
  \chi(\mathcal O_{K3})\,\chi(\mathcal O_E)=2\cdot0=0,
\]
and the compact chiral Hodge-supertrace invariant is
\[
  \kappa_{\mathrm{ch}}(K3\times E)
  =
  h^{0,0}-h^{0,1}+h^{0,2}-h^{0,3}
  =
  1-1+1-1=0.
\]
The Heisenberg--Mukai specialisation supplies a different chiral
invariant,
\[
  \kappa_{\mathrm{ch}}^{\mathrm{Heis}}(K3\times E)=3.
\]
The BKM and fibre values are
\[
  \kappa_{\mathrm{BKM}}(\Delta_5)=5,\qquad
  \kappa_{\mathrm{fiber}}(K3)=\mathrm{rk}\,\widetilde H(K3,\mathbb Z)
  =24.
\]
Thus the spectrum \(\{0,3,5,24\}\) records four constructions, with the
compact value \(0\) appearing simultaneously in \(\kappa_{\mathrm{cat}}\)
and \(\kappa_{\mathrm{ch}}\).

For a CY$_3$ input, the specialised chiral output of $\Phi_3$ is native
$\Eone$ unless a centre, double, or representation-category construction
is added. The identity in this note is a rank-$48$ Heisenberg residue
calculation on the \(K3\times E\) branch. The local toric statement is
\(\mathrm{CoHA}(\mathbb C^3)=Y^+\); a
\(\mathcal W_{1+\infty}\) comparison requires double or centre data.

\paragraph{Proof obligations.}
The finite coefficient theorem is
\[
  q^2\bigl[(2\pi i z)^2\Delta_5^{-2}\bigr]_{H_1,z=0}
  =
  \chi_{\Heis}^{\mathrm{Fock}}(q)
  \quad \text{through }q^{24},
\]
with the displayed coefficients computed by two independent finite
series formulae and the truncated Euler product, three algorithms in
total. The exact all-degree product identity requires the
Gritsenko--Nikulin separating degeneration and the standard rank-$48$
Heisenberg character formula. Any stronger
comparison must specify: primitive \(\Delta_5\) or monic \(\Dfive\);
unnormalised \(\PhiUn\) or OP scalar \(\PhiOP\); CHL, Nikulin, or
Gritsenko--Clery indexing; compact \(\kappa_{\mathrm{ch}}\) or
Heisenberg \(\kappa_{\mathrm{ch}}^{\mathrm{Heis}}\); Heisenberg Fock,
Virasoro, class-\(\mathcal S\), OP/DT, or interior Siegel-modular object.
